\section{Experiments}

We design four experiments aligned to RQ1--RQ4. The formal definitions
of behavioral uncertainty, the six pressure levels, the flip metric, and
the external calibrator are given in \autoref{sec:methods}; here we
describe the analyses applied to those measurements.

\paragraph{Experiment 1: Uncertainty and susceptibility (RQ1).}
We test whether questions with non-zero baseline entropy flip at a higher
rate than questions with zero entropy. Using the per-question Shannon
entropy $H(q;\mathcal{M})$ over $K = 25$ pre-pressure samples, we split
correct-at-$T_0$ questions into a \emph{certain} group ($H = 0$, all 25
samples agree) and an \emph{uncertain} group ($H < 0$, samples spread
across at least two options). For each model we compute the per-question
flip rate across pressure turns $T_1$--$T_6$ within each group and
compare the two groups with a Mann-Whitney $U$ test, reporting Wilson
95\% confidence intervals on each group's flip rate.

\paragraph{Experiment 2: Pressure-uncertainty interaction (RQ2).}
We test whether escalating social pressure increases the likelihood of
sycophantic flipping, and whether this effect depends on the model's
pre-pressure uncertainty (i.e., whether the marginal effect of each
additional pressure turn varies with the baseline entropy $H(q)$).

Concretely, we fit a logistic regression predicting the per-turn flip
outcome from pressure level $P$ (the turn index, defined in
\autoref{tab:pressure_levels}), baseline entropy $H$, and their
interaction $P \times H$:
\[
P(\mathrm{flip}) =
\sigma\!\left(\beta_0 + \beta_1 P + \beta_2 H + \beta_3 P H\right).
\]
A positive interaction coefficient $\beta_3 > 0$ would indicate that the
marginal effect of each additional pressure turn is larger when the
model is already uncertain: pressure amplifies pre-existing uncertainty
rather than acting independently of it. A coefficient near zero
would indicate that pressure has a uniform effect regardless of baseline
uncertainty, and a negative coefficient would indicate that pressure is
\emph{less} effective on uncertain questions (a pattern we do not
observe).

\paragraph{Experiment 3: Per-model trajectories (RQ3).}
We test whether the uncertainty-conditioned pressure effect holds
uniformly across model families and capability tiers. We fit a
mixed-effects logistic regression with the same fixed effects as
Experiment~2 and a model-level random intercept, which lets the
intercept (overall susceptibility) shift per model while estimating a
shared pressure-entropy interaction. We additionally compute per-model
flip rates at every turn $T_1$--$T_6$ and fit a linear trend
(percentage points per turn) within each model. The sign of the
per-model slope partitions models into two qualitative groups, which we
analyze in \autoref{sec:results}.

\paragraph{Experiment 4: Reasoning-trace capitulation (RQ4).}
We examine whether capitulation is detectable in the chain-of-thought
trace before the final answer changes. Using the free-form CoT prompt
and external calibrator described in
\autoref{sec:methods:calibrator}, we obtain per-step beliefs
$\hat{c}_\ell$ and confidences $\hat{\gamma}_\ell$ for every reasoning
trace in every dialogue. From these we compute, at each pressure turn
$t$: (i) the belief entropy over the run-level majority beliefs, which
measures how much the calibrator's read of the reasoning trace
fluctuates across repeated samples; (ii) the majority-belief
correctness, i.e.\ the fraction of runs whose final calibrated belief
matches $a^*$; and (iii) the mean calibrated confidence at the final
reasoning step, from which we derive the overconfidence gap
(Conf.\ $-$ Maj.\ Correct). Together these three quantities let us ask
whether the trace exposes capitulation earlier than the final answer
trajectory and whether confidence tracks the underlying belief shift.
